\documentclass[
  12pt,
  oneside,
  a4paper,
  english,
  brazilian
  ]{abntex2}

% ---
% Pacotes básicos
% ---
\usepackage[scaled]{helvet}
\renewcommand{\familydefault}{\sfdefault}
\usepackage[T1]{fontenc}
\usepackage[utf8]{inputenc}
\usepackage{indentfirst}
\usepackage{color}
\usepackage{graphicx}
\usepackage{microtype}
\usepackage{lipsum}
\usepackage{amsthm}
\usepackage{url}
\setlength{\marginparwidth}{2cm}
\usepackage{todonotes}
\usepackage{tikz}
\usetikzlibrary{arrows.meta, positioning}
\usepackage{amsmath, amssymb}
\usepackage{enumitem}

% ---
% Pacotes de citações
% ---
\usepackage[alf]{abntex2cite}

% ---
% Ajustes de fonte para títulos
% ---
\renewcommand{\ABNTEXchapterfont}{\sffamily\bfseries}
\renewcommand{\ABNTEXchapterfontsize}{\LARGE}
\renewcommand{\ABNTEXsectionfont}{\sffamily\bfseries}
\renewcommand{\ABNTEXsectionfontsize}{\Large}
\renewcommand{\ABNTEXsubsectionfont}{\sffamily\bfseries}
\renewcommand{\ABNTEXsubsectionfontsize}{\large}

\renewcommand{\thepage}{\arabic{page}}
\renewcommand{\thechapter}{\arabic{chapter}}
\renewcommand{\thesection}{\thechapter.\arabic{section}}
\renewcommand{\thesubsection}{\thechapter.\arabic{subsection}}
\renewcommand{\thefigure}{\arabic{figure}}
\renewcommand{\thetable}{\arabic{table}}

\captionstyle{\sffamily\bfseries}
\captionnamefont{\sffamily\bfseries}
\captiontitlefont{\sffamily}

% ---
% Ambientes de teorema
% ---
\theoremstyle{definition}
\newtheorem{definicao}{Definição}[chapter]
\newtheorem{teorema}{Teorema}[chapter]
\newtheorem{lema}{Lema}[chapter]
\newtheorem{exemplo}{Exemplo}[chapter]
\newtheorem{propriedade}{Propriedade}[chapter]

% ---
% Proteção para Metadados do PDF
% ---
\pdfstringdefDisableCommands{%
  \def\\{ }%
  \def\par{ }%
  \def\vspace#1{ }%
  \def\sffamily{}%
  \def\bfseries{}%
}

% ---
% Arquivo de Bibliografia Embutido
% ---
\begin{filecontents*}{referencias.bib}

@misc{fabricio2019,
  author = {Siqueira, Fabrício and others},
  title  = {Portal OBMEP: Material Teórico -- Tópicos Adicionais --
            Introdução à Teoria Dos Grafos -- Parte 1},
  year   = {2019}
}

@book{boaventura2009,
  author    = {Boaventura, Oswaldo and Jurkiewicz, Samuel},
  title     = {Grafos},
  publisher = {Editora Blucher},
  year      = {2009}
}

@misc{silva2020,
  author = {Da Silva, Luiz Fernando},
  title  = {Uma Proposta de Sequência Didática para o Ensino Médio:
            Introdução ao Estudo da Teoria dos Grafos com GeoGebra},
  year   = {2020},
  note   = {6 mai. 2020}
}

@book{cormen2009,
  author    = {Cormen, T. H. and Leiserson, C. E. and Rivest, R. L. and Stein, C.},
  title     = {Introduction to Algorithms},
  edition   = {3},
  publisher = {MIT Press},
  address   = {Cambridge, Mass.},
  year      = {2009}
}

@mastersthesis{custodio2025,
  author  = {Custodio, Diego Scanavachi},
  title   = {Um estudo sobre grafos e sua aplicabilidade no ensino médio},
  school  = {Universidade de São Paulo},
  address = {São Carlos},
  year    = {2025},
  note    = {Disponível em:
             \url{https://www.teses.usp.br/teses/disponiveis/55/55136/tde-24072025-110655/}.
             Acesso em: 4 nov. 2025}
}

@manual{brasil2002,
  organization = {Brasil},
  title        = {{PCN+} Ensino Médio: Orientações Educacionais Complementares
                  aos Parâmetros Curriculares Nacionais -- Ciências da Natureza,
                  Matemática e suas Tecnologias},
  address      = {Brasília},
  publisher    = {MEC/SEMTEC},
  year         = {2002}
}

@misc{figueredo,
  author = {Figueredo, J.},
  title  = {Teoria dos Grafos: Representação de Grafos},
  year   = {s.d.}
}

@misc{jurkiewicz2007,
  author = {Jurkiewicz, S.},
  title  = {Grafos: Uma introdução},
  year   = {2007},
  note   = {Programa de Iniciação Científica da OBMEP 2007}
}

@book{lopes2010,
  author    = {Lopes, M. L. M. L.},
  title     = {Grafos: jogos e desafios},
  address   = {Rio de Janeiro},
  publisher = {IM/UFRJ},
  year      = {2010}
}

@mastersthesis{malta2008,
  author  = {Malta, G. H. S.},
  title   = {Grafos no ensino médio: uma inserção possível},
  school  = {Universidade Federal do Rio Grande do Sul},
  address = {Porto Alegre},
  year    = {2008}
}

@book{santos2007,
  author    = {Santos, J. P. O. and Mello, M. P. and Murari, I. T. C.},
  title     = {Introdução à Análise Combinatória},
  edition   = {4},
  address   = {Rio de Janeiro},
  publisher = {Editora Ciência Moderna},
  year      = {2007}
}

@article{rodrigues2021,
  author  = {Rodrigues, A. B. and Silva, C. D.},
  title   = {O ensino de grafos e a cultura digital: novas perspectivas
             para a matemática escolar},
  journal = {Revista de Educação Matemática},
  volume  = {10},
  number  = {2},
  pages   = {1--15},
  year    = {2021}
}

@book{diestel2017,
  author    = {Diestel, Reinhard},
  title     = {Graph Theory},
  edition   = {5},
  publisher = {Springer},
  address   = {Berlin},
  year      = {2017}
}

@book{wilson1996,
  author    = {Wilson, Robin J.},
  title     = {Introduction to Graph Theory},
  edition   = {4},
  publisher = {Longman},
  address   = {London},
  year      = {1996}
}

@book{bondy1976,
  author    = {Bondy, J. A. and Murty, U. S. R.},
  title     = {Graph Theory with Applications},
  publisher = {North-Holland},
  address   = {New York},
  year      = {1976}
}

@book{rosen2012,
  author    = {Rosen, Kenneth H.},
  title     = {Discrete Mathematics and Its Applications},
  edition   = {7},
  publisher = {McGraw-Hill},
  address   = {New York},
  year      = {2012}
}

@article{euler1736,
  author  = {Euler, Leonhard},
  title   = {Solutio problematis ad geometriam situs pertinentis},
  journal = {Commentarii Academiae Scientiarum Imperialis Petropolitanae},
  volume  = {8},
  pages   = {128--140},
  year    = {1741},
  note    = {Apresentado em 1736}
}

@misc{brasil1996,
  author   = {Brasil},
  title    = {Lei n\textordmasculine{} 9.394, de 20 de dezembro de 1996:
              Estabelece as diretrizes e bases da educação nacional},
  address  = {Brasília},
  year     = {1996},
  note     = {Diário Oficial da União, Brasília, DF, 23 dez. 1996.
              Disponível em: \url{http://www.planalto.gov.br/ccivil_03/leis/l9394.htm}.
              Acesso em: 10 jan. 2026}
}

\end{filecontents*}

% ---
% Informações de dados para CAPA e FOLHA DE ROSTO
% ---
\titulo{PROPOSTA DE INSERÇÃO DO ENSINO DE GRAFOS: PARA ALUNOS DO ENSINO MÉDIO DA REDE PÚBLICA DE ENSINO}
\autor{JOSÉ GUILHERME DANTAS SILVA}
\local{Ilhéus-BA}
\data{2026}
\orientador{Prof. Dr. Romenique da Rocha Silva}
\coorientador[Coorientadora:]{Prof\textsuperscript{a}. Ma. Sara Meira Moutta}
\instituicao{%
  UNIVERSIDADE ESTADUAL DE SANTA CRUZ
  \par
  Departamento de Ciências Exatas
  \par
  Licenciatura em Matemática}
\tipotrabalho{Trabalho de Conclusão de Curso}
\preambulo{Trabalho de Conclusão de Curso apresentado ao curso de Licenciatura em Matemática, Departamento de Ciências Exatas, da UNIVERSIDADE ESTADUAL DE SANTA CRUZ, como requisito parcial para a obtenção do grau de Licenciado em Matemática.}

% ---
% Redefinição da Capa
% ---
\renewcommand{\imprimircapa}{%
  \begin{capa}%
    \centering
    \includegraphics[width=3cm]{Brasão-UESC---png1.png}\\[0.4cm]
    {\ABNTEXchapterfont\large\imprimirinstituicao}

    \vspace{2.5cm}

    {\ABNTEXchapterfont\large\imprimirautor}

    \vspace*{\fill}

    {\ABNTEXchapterfont\bfseries\LARGE\imprimirtitulo}

    \vspace*{\fill}

    {\large\imprimirlocal}\\[0.3cm]
    {\large\imprimirdata}\\[1cm]
  \end{capa}
}

% ---
% Configurações de aparência do PDF final
% ---
\definecolor{blue}{RGB}{41,5,195}

\hypersetup{
    pdftitle  = {\imprimirtitulo},
    pdfauthor = {\imprimirautor},
    colorlinks = true,
    linkcolor  = blue,
    citecolor  = blue,
    urlcolor   = blue
}

\setlength{\parindent}{1.3cm}
\setlength{\parskip}{0.2cm}

% ============================================================
% INÍCIO DO DOCUMENTO
% ============================================================
\begin{document}
\selectlanguage{brazilian}
\frenchspacing

% ----------------------------------------------------------
% ELEMENTOS PRÉ-TEXTUAIS
% ----------------------------------------------------------
\pretextual

\imprimircapa
\imprimirfolhaderosto*

% ---
% Folha de aprovação
% ---
\begin{folhadeaprovacao}
  \begin{center}
    {\ABNTEXchapterfont\large\imprimirautor}
    \vspace*{\fill}\vspace*{\fill}
    \begin{center}
      \ABNTEXchapterfont\bfseries\Large\imprimirtitulo
    \end{center}
    \vspace*{\fill}
    \hspace{.45\textwidth}
    \begin{minipage}{.5\textwidth}
        \imprimirpreambulo
    \end{minipage}%
    \vspace*{\fill}
  \end{center}

  Aprovado em: Ilhéus, \_\_ de \_\_\_\_\_\_\_\_\_\_\_\_ de \_\_\_\_.
  \begin{center}
    \vspace*{0.5cm}
    \textbf{BANCA EXAMINADORA} \\
    \vspace*{1.5cm}
    \hrulefill \\
    \textbf{\imprimirorientador} (Orientador) \\
    Universidade Estadual de Santa Cruz \\
    \vspace*{1.5cm}
    \hrulefill \\
    \textbf{\imprimircoorientador} (Coorientadora) \\
    Universidade Estadual de Santa Cruz \\
    \vspace*{1.5cm}
    \hrulefill \\
    \textbf{Prof\textsuperscript{a}. Dr\textsuperscript{a}.\ Jurema Lindote Botelho Peixoto} (Examinadora) \\
    Universidade Estadual de Santa Cruz \\
  \end{center}
\end{folhadeaprovacao}

% ---
% Dedicatória
% ---
\begin{dedicatoria}
   \vspace*{\fill}
   \centering
   \noindent
   \textit{Dedico este trabalho in memoriam a meu pai, Geraldo Bezerra da Silva,
           como homenagem à sua vida e legado.}
   \vspace*{\fill}
\end{dedicatoria}

% ---
% Agradecimentos
% ---
\begin{agradecimentos} \todo[inline, color=yellow]{AJUSTAR A PARTE DE AGRADECIMENTOS}
Gostaria de começar agradecendo a Deus, que, com sua infinita bondade e
misericórdia, me deu forças e me capacitou para chegar até aqui. Sei, inclusive,
que Ele ainda me fará ir muito mais longe.

Agradeço à Universidade Estadual de Santa Cruz (UESC), instituição de excelência
que me proporcionou anos de grande enriquecimento, tanto intelectual quanto
espiritual. Tenho certeza de que retornarei a esta casa para cursar uma
pós-graduação e, futuramente, contribuir como docente.

Registro minha gratidão, de forma especial, ao meu orientador, professor
Dr.\ Romenique da Rocha Silva, cuja linha de pesquisa, especialmente no doutorado,
despertou em mim ainda mais fascínio pela área de computação e grafos. Esse
encantamento motivou a realização deste trabalho, com o qual pretendo contribuir
para a educação básica. Agradeço por sua presença constante, dedicação e empenho
em ler esta monografia com atenção, tornando-a uma pesquisa de excelência. Sou
igualmente grato pela amizade que cultivamos ao longo desse tempo e sei que ainda
estaremos juntos em outros projetos.

Expresso também meus agradecimentos a duas instituições da educação básica: o
Colégio Municipal Antônio Lima de Oliveira e o Colégio Estadual do Salobrinho.
Obrigado por acreditarem em mim, mesmo antes da hora, e por me permitirem
realizar o sonho de ser professor.

A todos os alunos que passaram por mim, deixo
o reconhecimento: vocês foram fundamentais para que eu continuasse. Em muitos
momentos, pensei em desistir, mas vocês me mostraram a beleza da docência e me
deram forças para seguir.

Sou profundamente grato às pessoas que tornaram este sonho possível. À minha mãe,
Edileusa Gualberto Dantas (Inha), à minha avó, Doraci Maria da Silva (Vó Peró),
e às minhas irmãs, Victória Regina (Vick) e Claudineia (Lela), que nunca
deixaram de acreditar em mim, mesmo quando eu duvidava de mim mesmo.

Ao meu gato Bimbo, que esteve comigo em meio a noites em claro, compartilhando
angústias, com fórmulas espalhadas em papéis na garagem.

Ao meu colega e amigo Juan Cerqueira (Calopsita), meu parceiro de graduação que
se tornou um irmão: sem o seu apoio, companheirismo e cumplicidade, eu não teria
resistido. Sou imensamente grato por sua vida e sei que ainda vamos brilhar muito
juntos.

Por fim, agradeço aos meus supervisores de estágio e PIBID, Prof.\ Wellingthon
Fraga, Prof\textsuperscript{a}.\ Luana Cerqueira e Prof\textsuperscript{a}.\
Daniella Feitosa. Só tenho a agradecer a vocês por todo o ensinamento docente e
pela oportunidade de vivenciar experiências em sala de aula. Não poderia deixar
de mencionar também meu professor do ensino médio, Ms.\ Edmilson Chaves, que
despertou em mim, ainda na educação básica, o desejo de me tornar docente, e a
professora Dr\textsuperscript{a}.\ Jurema Peixoto, que acreditou em mim desde o
início da graduação.
\end{agradecimentos}

% ---
% Epígrafe
% ---
\begin{epigrafe}
    \vspace*{\fill}
  \begin{flushright}
    \textit{``A vida é como a matemática, criamos raízes \\
    que o mundo tenta subtrair de nós, porém o que é somado \\
    de bom nos ajudará a ter a solução exata do que precisamos.''\\
    (Ester Menezes)}
  \end{flushright}
\end{epigrafe}

% ---
% Resumos
% ---
\begin{resumo}
 [Insira aqui o texto do seu resumo de 150 a 500 palavras]
 \vspace{\onelineskip}
 \noindent
 \textbf{Palavras-chave}: Teoria dos Grafos. Ensino Médio. Matemática. BNCC.
\end{resumo}

\begin{resumo}[Abstract]
 \begin{otherlanguage*}{english}
   [Insert the abstract text here]
   \vspace{\onelineskip}
   \noindent
   \textbf{Keywords}: Graph Theory. High School. Mathematics. BNCC.
 \end{otherlanguage*}
\end{resumo}

\pdfbookmark[0]{\listfigurename}{lof}
\listoffigures*
\cleardoublepage

\begin{siglas}
  \item[BNCC] Base Nacional Comum Curricular
  \item[LDB]  Lei de Diretrizes e Bases da Educação Nacional
  \item[UESC] Universidade Estadual de Santa Cruz
\end{siglas}

\pdfbookmark[0]{\contentsname}{toc}
\tableofcontents*
\cleardoublepage

% ----------------------------------------------------------
% ELEMENTOS TEXTUAIS
% ----------------------------------------------------------
\textual

% ==========================================================
\chapter{INTRODUÇÃO}
% ==========================================================

A contextualização do ensino da Matemática na Educação Básica, particularmente
no Ensino Médio, demanda uma constante atualização curricular que dialogue com a
complexidade e a interconectividade do mundo contemporâneo. Nesse cenário, a
Teoria dos Grafos destaca-se não apenas como um ramo fundamental da Matemática
Discreta, mas como uma ferramenta indispensável para a modelagem e a compreensão
de fenômenos reais.

A relevância desse campo transcende a abstração teórica, manifestando-se
diretamente em estruturas que definem a sociedade atual, tais como a arquitetura
de redes sociais, o planejamento eficiente de sistemas de transporte e logística,
o funcionamento de algoritmos de busca e a análise de grandes volumes de dados
(\textit{Big Data}). Compreender grafos é, portanto, compreender a linguagem
estrutural por trás dos problemas de otimização e conectividade que permeiam o
cotidiano.

Para além de sua onipresença tecnológica, o estudo de grafos apresenta uma vasta
aplicabilidade prática em áreas estratégicas do conhecimento, servindo como base
para avanços na Ciência da Computação, Engenharia, Biologia (na análise de redes
metabólicas e genéticas) e Economia. Do ponto de vista pedagógico, a introdução
desses conceitos no ambiente escolar contribui significativamente para o
desenvolvimento de competências cognitivas fundamentais. O trabalho com vértices
e arestas estimula o pensamento relacional e a capacidade de abstração,
ferramentas essenciais para a modelagem matemática, a resolução de problemas não
lineares e a interpretação de estruturas complexas.

A pertinência dessa abordagem é amplamente discutida na literatura acadêmica
atual. \citeonline{rodrigues2021}, por exemplo, analisam como a inserção de
tópicos de Teoria dos Grafos pode favorecer uma formação matemática mais robusta
e conectada ao universo digital, alinhando o currículo escolar às demandas da
contemporaneidade. Sob essa ótica, defende-se que a transposição didática desse
conteúdo para o Ensino Médio é uma possibilidade real e metodologicamente
fundamentada.

A Teoria dos Grafos, área da Matemática Discreta, dedica-se ao estudo dessas
estruturas formadas por vértices e arestas, possibilitando a representação e a
análise de relações e interações em diferentes contextos. Sua presença é
identificada em diversos campos, como sistemas de transporte público, redes
sociais, circuitos elétricos, algoritmos de busca na internet, logística
empresarial e até em áreas das ciências biológicas, como no estudo das cadeias
alimentares.

Apesar da ampla aplicabilidade, a Teoria dos Grafos ainda não recebe a devida
atenção no Ensino Médio da rede pública, cujos currículos e livros didáticos
priorizam conteúdos considerados clássicos da Matemática, muitas vezes
desvinculados das demandas tecnológicas e sociais contemporâneas. Essa lacuna
dificulta que os estudantes compreendam o potencial da Matemática para analisar e
propor soluções para problemas atuais.

No contexto do Novo Ensino Médio e da Base Nacional Comum Curricular (BNCC), que
enfatizam a contextualização, a interdisciplinaridade e o desenvolvimento do
pensamento computacional, a ausência da Teoria dos Grafos mostra-se ainda mais
significativa. Nesse sentido, a questão central que orienta este trabalho consiste
em investigar de que maneira é possível propor a inserção teórica da Teoria dos
Grafos no Ensino Médio da rede pública, de forma contextualizada, interdisciplinar
e em consonância com as diretrizes da BNCC.

Para dar conta dessa questão, o presente trabalho está organizado em capítulos
que se articulam entre si. O Capítulo~\ref{cap:fundamentos} apresenta os
fundamentos teóricos da Teoria dos Grafos e o estado da arte sobre seu ensino.
O Capítulo~\ref{cap:bncc} realiza uma análise documental da BNCC, identificando
as conexões entre a Teoria dos Grafos e as competências e habilidades previstas
para o Ensino Médio. O Capítulo~\ref{cap:proposta}, por fim, apresenta a proposta
didática propriamente dita, articulando os conceitos teóricos e as atividades
sugeridas com as diretrizes curriculares e as estratégias pedagógicas
recomendadas pela literatura.

\section{Objetivo Geral}

Elaborar uma abordagem teórica para a inserção da Teoria dos Grafos na Matemática
do Ensino Médio da rede pública de ensino, destacando possibilidades de
contextualização, interdisciplinaridade e inovação curricular no ensino da
Matemática.

\section{Objetivos Específicos}

\begin{itemize}
  \item Investigar a relevância da Teoria dos Grafos no contexto da Educação do
        Ensino Médio e sua relação com as competências e habilidades previstas na
        BNCC;
  \item Identificar exemplos de aplicações de grafos em situações do cotidiano e
        de outras áreas do conhecimento, capazes de tornar o aprendizado mais
        significativo para os estudantes;
  \item Pesquisar potenciais estratégias didáticas para a introdução dos conceitos
        de grafos no Ensino Médio, com foco em abordagens interdisciplinares e
        contextualizadas;
  \item Analisar os principais desafios e limitações para a implementação desse
        conteúdo na rede pública de ensino;
  \item Contribuir para o debate acadêmico sobre a modernização do ensino de
        Matemática, oferecendo uma proposta teórica fundamentada que possa
        subsidiar futuras pesquisas e práticas de sala de aula.
\end{itemize}

% ==========================================================
\chapter{METODOLOGIA}
% ==========================================================

Este trabalho caracteriza-se como uma pesquisa de natureza qualitativa, de
caráter bibliográfico e exploratório, voltada para a elaboração de uma proposta
de inserção da Teoria dos Grafos no Ensino Médio da rede pública de ensino.

A escolha por essa abordagem metodológica justifica-se pela natureza do problema
investigado, que busca compreender lacunas e possibilidades no currículo escolar
sem prever aplicação prática em sala de aula, mas sim a construção de um
referencial teórico e de sugestões didáticas fundamentadas em pesquisas já
existentes.

Inicialmente, foi realizado um levantamento bibliográfico em livros, artigos
científicos, dissertações e teses que tratam do ensino de Matemática na Educação
Básica, da Matemática Discreta com ênfase na Teoria dos Grafos. Também foram
analisados estudos sobre inovação curricular e metodologias interdisciplinares
que utilizam grafos como recurso pedagógico.

Além disso, foi realizada uma análise documental da Base Nacional Comum
Curricular (BNCC) e de currículos estaduais e municipais, a fim de identificar as
competências e habilidades previstas para o Ensino Médio que se articulam com a
Teoria dos Grafos.

A partir desse referencial bibliográfico e da análise documental, o trabalho
foi estruturado em três eixos complementares: (i) a apresentação dos
fundamentos teóricos e do estado da arte sobre o ensino de grafos,
reunidos no Capítulo~\ref{cap:fundamentos}; (ii) a análise documental da BNCC
e suas implicações para o ensino de grafos, desenvolvida no
Capítulo~\ref{cap:bncc}; e (iii) a proposta didática propriamente dita, que
articula os dois eixos anteriores com estratégias pedagógicas e atividades
contextualizadas, exposta no Capítulo~\ref{cap:proposta}. Essa organização visa
garantir coerência interna entre o rigor matemático, a dimensão didática e o
enquadramento normativo da proposta.

Por fim, foi realizada uma discussão crítica dos potenciais benefícios da proposta
para o desenvolvimento do raciocínio lógico, da capacidade de resolução de
problemas e da interdisciplinaridade, bem como das dificuldades e limitações que
poderiam emergir em um cenário de eventual implementação, como a necessidade de
formação docente específica, a sobrecarga curricular e a escassez de recursos nas
escolas públicas.

% ==========================================================
\chapter{FUNDAMENTOS TEÓRICOS E ESTADO DA ARTE}
\label{cap:fundamentos}
% ==========================================================

Este capítulo apresenta os fundamentos matemáticos da Teoria dos Grafos e uma
revisão da literatura sobre seu ensino na Educação Básica. A exposição segue uma
progressão natural dos conceitos --- do mais elementar ao mais elaborado ---
permitindo que o leitor compreenda a estrutura lógica da teoria antes de apreciá-
la em sua dimensão didática e aplicada. A abordagem adota a tradição consolidada
na literatura especializada \cite{diestel2017, wilson1996, boaventura2009}.

% ----------------------------------------------------------
\section{Breve Histórico e o Problema de Königsberg}
% ----------------------------------------------------------

A Teoria dos Grafos é considerada um dos campos mais antigos da Matemática
Discreta. Seu surgimento é atribuído ao matemático suíço Leonhard Euler, que em
1736 publicou a solução do famoso \textit{Problema das Sete Pontes de
Königsberg} \cite{euler1736}. O problema consistia em determinar se era possível
percorrer todas as sete pontes da cidade de Königsberg (atual Kaliningrado,
Rússia), que ligavam as margens do Rio Pregel a duas ilhas, passando por cada
ponte exatamente uma vez e retornando ao ponto de partida.

Euler demonstrou que tal percurso era impossível e, ao fazê-lo, formalizou os conceitos de vértice e
aresta, lançando as bases da Teoria dos Grafos \cite{boaventura2009}. Essa história é, por si só, um ponto de partida pedagógico valioso: ela mostra
aos estudantes que a Matemática nasce da tentativa de resolver problemas
concretos, reforçando a dimensão aplicada do conhecimento matemático
\cite{malta2008}. O problema das pontes será retomado adiante quando
discutirmos os conceitos de caminho e circuito euleriano (Seção~\ref{sec:euleriano}).

% ----------------------------------------------------------
\section{Definições Matemáticas Fundamentais}
% ----------------------------------------------------------

% ----------------------------------------------------------
\subsection{Grafo}
% ----------------------------------------------------------

\begin{definicao}[Grafo]
\label{def:grafo}
  Um \textit{grafo} $\mathcal{G}$ é um par ordenado
  $\mathcal{G} = (\mathcal{V},\, \mathcal{A})$, em que $\mathcal{V}$ é um
  conjunto não vazio de elementos denominados \textit{vértices} (ou nós) e
  $\mathcal{A}$ é um conjunto de pares não ordenados de elementos de
  $\mathcal{V}$, denominados \textit{arestas} \cite[p.~1]{boaventura2009}.
\end{definicao}

De maneira intuitiva, pode-se imaginar um grafo como um diagrama composto por
pontos (os vértices) conectados por linhas (as arestas). Quando dois vértices
estão ligados por uma aresta, diz-se que eles são \textit{adjacentes}
\cite{rosen2012}. A quantidade de vértices de $\mathcal{G}$ é chamada de
\textit{ordem} do grafo, enquanto a quantidade de arestas é chamada de
\textit{tamanho}.

% ----------------------------------------------------------
\subsection{Grau de um Vértice e o Lema do Aperto de Mão}
% ----------------------------------------------------------

\begin{definicao}[Grau de um vértice]
\label{def:grau}
  O \textit{grau} de um vértice $v$, denotado por $\delta(v)$, corresponde ao
  número de arestas que incidem sobre ele \cite{bondy1976}. Um vértice de grau
  zero é denominado \textit{vértice isolado}.
\end{definicao}

\begin{lema}[Lema do Aperto de Mão]
\label{lema:handshaking}
  Em qualquer grafo $\mathcal{G} = (\mathcal{V},\, \mathcal{A})$, a soma dos
  graus de todos os vértices é igual ao dobro do número de arestas:
  \begin{equation}
    \sum_{v \in \mathcal{V}} \delta(v) = 2\,|\mathcal{A}|.
\label{eq:lema_aperto}
  \end{equation}
\end{lema}

\begin{proof}
  Cada aresta $\{u, v\} \in \mathcal{A}$ contribui com $+1$ ao grau de $u$ e
  com $+1$ ao grau de $v$. Logo, ao somarmos os graus de todos os vértices,
  cada aresta é contabilizada exatamente duas vezes, o que prova a
  igualdade~\eqref{eq:lema_aperto}.
\end{proof}

\begin{propriedade}
  Em qualquer grafo, o número de vértices de grau ímpar é sempre par.
\end{propriedade}

\begin{proof}
  Seja $I$ o conjunto dos vértices de grau ímpar e $P$ o conjunto dos vértices
  de grau par. Pelo Lema~\ref{lema:handshaking},
  \[
    \sum_{v \in I} \delta(v) + \sum_{v \in P} \delta(v) = 2\,|\mathcal{A}|.
  \]
  O segundo somatório e o lado direito são pares; logo, $\sum_{v \in I}
  \delta(v)$ também é par. Como cada parcela é ímpar, o número de parcelas
  --- isto é, $|I|$ --- deve ser par.
\end{proof}

% ----------------------------------------------------------
\subsection{Grafos Não Dirigidos e Grafos Dirigidos}
% ----------------------------------------------------------

\begin{definicao}[Grafo não dirigido e dígrafo]
\label{def:dirigido}
  Um grafo é dito \textit{não dirigido} quando suas arestas não possuem
  orientação, ou seja, a relação entre dois vértices é simétrica. Um
  \textit{grafo dirigido} (ou \textit{dígrafo}) é aquele em que as arestas
  possuem uma orientação definida, representadas por \textit{arcos}; o arco
  $(u, v)$ parte de $u$ e chega em~$v$, sendo distinto do arco $(v, u)$
  \cite{cormen2009}.
\end{definicao}

Em situações práticas do cotidiano, grafos não dirigidos modelam relações
simétricas, como amizades em redes sociais, ao passo que dígrafos modelam
relações assimétricas, como o fluxo de tráfego em vias de mão única ou a
navegação entre páginas da internet \cite{rosen2012}.

% ----------------------------------------------------------
\subsection{Subgrafo e Grafo Completo}
% ----------------------------------------------------------

\begin{definicao}[Subgrafo]
  Dado um grafo $\mathcal{G} = (\mathcal{V},\, \mathcal{A})$, um
  \textit{subgrafo} $\mathcal{H}$ de $\mathcal{G}$ é um grafo tal que
  $\mathcal{V}(\mathcal{H}) \subseteq \mathcal{V}(\mathcal{G})$ e
  $\mathcal{A}(\mathcal{H}) \subseteq \mathcal{A}(\mathcal{G})$
  \cite{wilson1996}.
\end{definicao}

\begin{definicao}[Grafo completo]
\label{def:completo}
  Um grafo é dito \textit{completo} quando todo par de vértices distintos está
  ligado por exatamente uma aresta. O grafo completo com $n$ vértices é
  denotado por $K_n$ \cite{bondy1976}.
\end{definicao}

\begin{teorema}
  O grafo completo $K_n$ possui exatamente $\dfrac{n(n-1)}{2}$ arestas, e
  cada vértice possui grau $n - 1$.
\end{teorema}

\begin{proof}
  Cada um dos $n$ vértices é adjacente a todos os demais $n-1$ vértices. Assim,
  a soma dos graus é $n(n-1)$. Pelo Lema do Aperto de Mão
  (Lema~\ref{lema:handshaking}), o número de arestas é
  $\dfrac{n(n-1)}{2}$.
\end{proof}

% ----------------------------------------------------------
\subsection{Caminho, Ciclo e Conectividade}
% ----------------------------------------------------------

\begin{definicao}[Caminho]
  Um \textit{caminho} em um grafo é uma sequência de vértices distintos
  $v_1, v_2, \ldots, v_k$ tal que existe uma aresta entre cada par consecutivo
  $(v_i, v_{i+1})$ \cite{diestel2017}. O número de arestas percorridas é o
  \textit{comprimento} do caminho.
\end{definicao}

\begin{definicao}[Ciclo]
  Um \textit{ciclo} é um caminho fechado, isto é, um caminho em que o primeiro
  e o último vértice coincidem \cite{boaventura2009}. Um grafo sem ciclos é
  denominado \textit{acíclico} ou \textit{floresta}.
\end{definicao}

\begin{definicao}[Grafo conexo]
  Um grafo é denominado \textit{conexo} quando existe pelo menos um caminho
  entre qualquer par de vértices; caso contrário, é dito \textit{desconexo}
  \cite{rosen2012}. As partes maximais conexas de um grafo desconexo são
  chamadas de \textit{componentes conexas}.
\end{definicao}

% ----------------------------------------------------------
\subsection{Caminho e Circuito Euleriano}
\label{sec:euleriano}
% ----------------------------------------------------------

\begin{definicao}[Caminho euleriano e circuito euleriano]
\label{def:euleriano}
  Um \textit{caminho euleriano} é um caminho que percorre cada aresta do grafo
  exatamente uma vez. Um \textit{circuito euleriano} é um caminho euleriano
  que retorna ao vértice de origem \cite{euler1736}.
\end{definicao}

\begin{teorema}[Teorema de Euler]
\label{teo:euler}
  Um grafo conexo admite um circuito euleriano se, e somente se, todos os seus
  vértices possuem grau par. Admite um caminho euleriano (sem ser circuito) se,
  e somente se, possui exatamente dois vértices de grau ímpar
  \cite{diestel2017}.
\end{teorema}

% ----------------------------------------------------------
\section{O Ensino de Grafos na Literatura Acadêmica}
% ----------------------------------------------------------

Nesta seção, apresentamos uma breve revisão de trabalhos acadêmicos que discutiram a inserção da Teoria dos Grafos no Ensino Médio. A literatura mostra que, embora o tema ainda seja pouco explorado nos currículos oficiais, há um crescente interesse por parte de pesquisadores em Educação Matemática.

\citeonline{malta2008}, em sua dissertação de mestrado, compartilha da ideia de que ``a Teoria dos Grafos apresenta aspectos pertinentes que merecem espaço no currículo da Escola Básica''. A autora apoia-se no contexto histórico para considerar que pertencem à Introdução da Teoria dos Grafos os conceitos que surgiram até o século XVIII, desde \textit{O Problema das Sete Pontes de Königsberg} (1736) até o \textit{Icosian Game} (1856), desenvolvidos por Leonhard Euler e William Hamilton, respectivamente.

Como a Teoria dos Grafos foi formulada a partir de um problema, acreditamos que a metodologia de Resolução de Problemas pode ser utilizada para a abordagem de Grafos no Ensino Médio. Reafirmando a escolha desse caminho, as \textit{Orientações Educacionais Complementares aos Parâmetros Curriculares Nacionais} (PCN+) enfatizam que a resolução de problemas é peça central para o ensino de Matemática, pois o pensar e o fazer se mobilizam e se desenvolvem quando o indivíduo está engajado ativamente no enfrentamento de desafios \cite{brasil2002}. Essa competência não se desenvolve quando propomos apenas exercícios de aplicação dos conceitos e técnicas matemáticas, pois, neste caso, o que está em ação é uma simples transposição analógica \cite{brasil2002}.

Adicionalmente, a Lei de Diretrizes e Bases da Educação Nacional (LDB) preconiza que as metodologias de resolução de problemas sejam adequadas ao cotidiano discente \cite{brasil1996}. Tal diretriz busca fomentar o pensamento autônomo e a criação de soluções próprias, permitindo que o estudante desenvolva independência na capacidade de resolver problemas.

Nesse mesmo sentido, \citeonline{rodrigues2021} analisam como a inserção de tópicos de Teoria dos Grafos pode favorecer uma formação matemática mais robusta e conectada ao universo digital, alinhando o currículo escolar às demandas da contemporaneidade. Os autores destacam a importância do pensamento computacional e da modelagem de problemas do cotidiano.

Ainda sob a ótica da tecnologia, \citeonline{silva2020} propõe uma sequência didática para o ensino de grafos no Ensino Médio utilizando o software GeoGebra, evidenciando as potencialidades das ferramentas digitais para a visualização e exploração dos conceitos.

% ----------------------------------------------------------
\subsection{Exemplos Ilustrativos com Representação Gráfica}
% ----------------------------------------------------------

Os três exemplos a seguir ilustram as definições apresentadas e servirão de
referência para as atividades do Capítulo~\ref{cap:proposta}.

% ----------------------------------------------------------
\subsubsection{Exemplo 1 --- Grafo Não Dirigido}
% ----------------------------------------------------------

\begin{exemplo}
\label{ex:nao_dirigido}
Considere o grafo $\mathcal{G}_1 = (\mathcal{V}_1,\, \mathcal{A}_1)$ em que
\[
  \mathcal{V}_1 = \{A,\, B,\, C,\, D,\, E\}
  \quad \text{e} \quad
  \mathcal{A}_1 = \{AB,\, AC,\, AD,\, BC,\, CD,\, DE,\, EA\}.
\]

Esse grafo possui $|\mathcal{V}_1| = 5$ vértices e $|\mathcal{A}_1| = 7$
arestas. Aplicando o Lema do Aperto de Mão
(Lema~\ref{lema:handshaking}), a soma dos graus deve ser $2 \times 7 = 14$.
De fato: $\delta(A) = 4$, $\delta(B) = 2$, $\delta(C) = 3$, $\delta(D) = 3$,
$\delta(E) = 2$, totalizando $4 + 2 + 3 + 3 + 2 = 14$ \checkmark
\end{exemplo}

\begin{figure}[htb]
  \centering
  \begin{tikzpicture}[
      vertex/.style = {circle, draw, fill=blue!20,
                       minimum size=1.2cm, font=\sffamily\bfseries},
      edge/.style   = {thick}
  ]
    \node[vertex] (A) at (0,0)    {$A$};
    \node[vertex] (B) at (3,1)    {$B$};
    \node[vertex] (C) at (4,-2)   {$C$};
    \node[vertex] (D) at (1,-2.5) {$D$};
    \node[vertex] (E) at (-1,-1)  {$E$};

    \draw[edge] (A) -- (B);
    \draw[edge] (A) -- (C);
    \draw[edge] (A) -- (D);
    \draw[edge] (B) -- (C);
    \draw[edge] (C) -- (D);
    \draw[edge] (D) -- (E);
    \draw[edge] (E) -- (A);
  \end{tikzpicture}
  \caption{Grafo não dirigido $\mathcal{G}_1$ com 5 vértices e 7 arestas
           (Exemplo~\ref{ex:nao_dirigido}).}
  \label{fig:grafo_nao_dirigido}
\end{figure}

Esse grafo pode modelar, por exemplo, um mapa rodoviário simplificado em que os
vértices representam cidades e as arestas representam estradas de mão dupla que
as interligam \cite{jurkiewicz2007}. Verifique ainda que $\delta(A) = 4$ e
$\delta(B) = \delta(E) = 2$ são pares, enquanto $\delta(C) = \delta(D) = 3$
são ímpares; há exatamente dois vértices de grau ímpar, logo $\mathcal{G}_1$
admite caminho euleriano, mas não circuito euleriano
(Teorema~\ref{teo:euler}).

% ----------------------------------------------------------
\subsubsection{Exemplo 2 --- Grafo Dirigido (Dígrafo)}
% ----------------------------------------------------------

\begin{exemplo}
\label{ex:digrafo}
Considere o dígrafo $\mathcal{G}_2 = (\mathcal{V}_2,\, \mathcal{A}_2)$ em que
\[
  \mathcal{V}_2 = \{P,\, Q,\, R,\, S\}
  \quad \text{e} \quad
  \mathcal{A}_2 = \{(P,Q),\, (Q,R),\, (R,S),\, (S,P),\, (P,R)\}.
\]
Os arcos possuem direção: o arco $(P, Q)$ representa uma relação que parte de
$P$ e chega em~$Q$, mas não necessariamente o contrário.
\end{exemplo}

\begin{figure}[htb]
  \centering
  \begin{tikzpicture}[
      vertex/.style  = {circle, draw, fill=orange!25,
                        minimum size=1.2cm, font=\sffamily\bfseries},
      arc/.style     = {thick, -{Stealth[length=7pt]}},
      node distance  = 3cm
  ]
    \node[vertex] (P) at (0,0)   {$P$};
    \node[vertex] (Q) at (3,0)   {$Q$};
    \node[vertex] (R) at (3,-3)  {$R$};
    \node[vertex] (S) at (0,-3)  {$S$};
    \draw[arc] (P) -- (Q);
    \draw[arc] (Q) -- (R);
    \draw[arc] (R) -- (S);
    \draw[arc] (S) -- (P);
    \draw[arc] (P) -- (R);
  \end{tikzpicture}
  \caption{Dígrafo $\mathcal{G}_2$ com 4 vértices e 5 arcos orientados
           (Exemplo~\ref{ex:digrafo}).}
  \label{fig:digrafo}
\end{figure}

Essa estrutura é adequada para modelar vias de mão única em uma cidade ou o
fluxo de informações em uma rede de computadores \cite{cormen2009}. A distinção
entre grafos dirigidos e não dirigidos é explorada na
Atividade~\ref{ativ:redes}.

% ----------------------------------------------------------
\subsubsection{Exemplo 3 --- Grafo Completo $K_4$}
% ----------------------------------------------------------

\begin{exemplo}
\label{ex:k4}
O grafo completo $K_4$ possui $n = 4$ vértices e
$\dfrac{4 \times 3}{2} = 6$ arestas, sendo todo par de vértices adjacente.
Todos os vértices possuem grau $\delta = n - 1 = 3$ \cite{bondy1976}.
\end{exemplo}

\begin{figure}[htb]
  \centering
  \begin{tikzpicture}[
      vertex/.style = {circle, draw, fill=green!25,
                       minimum size=1.2cm, font=\sffamily\bfseries},
      edge/.style   = {thick}
  ]
    \node[vertex] (1) at (0, 2)   {$1$};
    \node[vertex] (2) at (2, 0)   {$2$};
    \node[vertex] (3) at (0,-2)   {$3$};
    \node[vertex] (4) at (-2, 0)  {$4$};

    \draw[edge] (1) -- (2);
    \draw[edge] (1) -- (3);
    \draw[edge] (1) -- (4);
    \draw[edge] (2) -- (3);
    \draw[edge] (2) -- (4);
    \draw[edge] (3) -- (4);
  \end{tikzpicture}
  \caption{Grafo completo $K_4$ com 4 vértices e 6 arestas
           (Exemplo~\ref{ex:k4}).}
  \label{fig:k4}
\end{figure}

Os grafos completos são frequentemente utilizados como casos-limite em problemas
de otimização combinatória, como o clássico Problema do Caixeiro Viajante, que
busca o menor caminho que passa por todos os vértices exatamente uma vez
\cite{rosen2012}. Esse problema é retomado na Atividade~\ref{ativ:caixeiro}.

% ==========================================================
\chapter{ANÁLISE DOCUMENTAL: A TEORIA DOS GRAFOS E A BNCC}
\label{cap:bncc}
% ==========================================================

Investigar a relevância da Teoria dos Grafos no contexto da Educação Básica e
sua relação com as competências e habilidades previstas na BNCC é essencial para
justificar sua inserção curricular. Este capítulo analisa a Base Nacional Comum
Curricular (BNCC) e documentos complementares, identificando pontos de
articulação com os conceitos de grafos apresentados no capítulo anterior.

% ----------------------------------------------------------
\section{A Matemática Discreta no Ensino Médio}
% ----------------------------------------------------------

A Matemática Discreta, da qual a Teoria dos Grafos é um dos ramos mais
proeminentes, tradicionalmente tem pouco espaço no currículo do Ensino Médio
brasileiro, que prioriza áreas como Álgebra, Geometria e Análise (Cálculo
diferencial e integral elementar). No entanto, a BNCC, ao propor a formação
integral do estudante e o desenvolvimento de competências para o século XXI,
abre brechas para a inclusão de tópicos de Matemática Discreta.

O documento oficial destaca a importância de o estudante ``identificar e
representar relações de interdependência entre grandezas, modelizar situações
diversas e resolver problemas que envolvem contagens, arranjos, combinações e
probabilidades'' (BNCC, p. 529). Embora mencione explicitamente análise
combinatória, a Teoria dos Grafos aparece como um desdobramento natural desses
estudos, fornecendo uma linguagem visual e relacional para problemas de
otimização e conectividade.

Além disso, as competências específicas de Matemática para o Ensino Médio, como
``utilizar estratégias, conceitos e procedimentos matemáticos para interpretar
situações em diversos contextos'' e ``propor ou participar de ações para
investigar desafios do mundo contemporâneo'', alinham-se perfeitamente à
abordagem por meio de grafos.

% ----------------------------------------------------------
\section{Relação com as Competências e Habilidades da BNCC}
% ----------------------------------------------------------

A Base Nacional Comum Curricular (BNCC) enfatiza a importância de desenvolver competências que preparem os alunos para enfrentar desafios do mundo real. Os conceitos formalizados no Capítulo~\ref{cap:fundamentos} e operacionalizados nas atividades do Capítulo~\ref{cap:proposta} articulam-se diretamente com diversas habilidades do documento. A Teoria dos Grafos oferece uma linguagem matemática poderosa que pode ser utilizada para resolver problemas complexos, como otimização de rotas e análise de redes. Podemos destacar as seguintes habilidades:

\begin{itemize}
  \item \textbf{EM13MAT101}: ``Interpretar situações econômicas, sociais e científicas por meio de gráficos, tabelas e medidas, identificando relações de dependência e representando-as por meio de modelos matemáticos''. A modelagem de redes sociais, rotas de transporte e problemas de conexão atende diretamente a essa habilidade.
  \item \textbf{EM13MAT201}: ``Propor e testar estratégias de resolução de problemas de contagem que envolvem agrupamentos ordenáveis ou não ordenáveis, de modo a construir e aplicar os conceitos de arranjo, combinação e permutação''. O Problema do Caixeiro Viajante e a contagem de rotas são exemplos concretos.
  \item \textbf{EM13MAT307}: ``Empregar diferentes estratégias para resolver problemas envolvendo a determinação de um percurso ou a organização de um circuito, considerando condições de adjacência, conectividade e otimização''. Essa habilidade é praticamente uma descrição do estudo de caminhos e circuitos eulerianos.
  \item \textbf{EM13MAT315}: ``Utilizar conhecimentos matemáticos para compreender e atuar no mundo digital, aplicando conceitos de algoritmos, pensamento computacional e modelagem''. A Teoria dos Grafos é fundamental para algoritmos de busca e sistemas de recomendação.
  \item \textbf{EM13MAT401}: ``Converter representações algébricas de funções polinomiais de 1º grau para representações geométricas no plano cartesiano, distinguindo os casos nos quais o comportamento é proporcional, recorrendo ou não a softwares ou aplicativos de álgebra e geometria dinâmica''. O estudo de grafos enriquece a visualização estrutural dessas relações e do pensamento analítico.
\end{itemize}

% ----------------------------------------------------------
\section{O Pensamento Computacional e a Interdisciplinaridade}
% ----------------------------------------------------------

A BNCC defende a promoção do pensamento algorítmico e computacional, e a Teoria dos Grafos oferece uma linguagem matemática que pode ser utilizada para resolver problemas complexos, como otimização de rotas e análise de redes. Ela fornece a base teórica ideal para ensinar esses conceitos sem a necessidade imediata de computadores, mas servindo de ponte estruturada para atividades tecnológicas mais avançadas.

Além disso, a interdisciplinaridade é um princípio estruturante do Novo Ensino Médio. A aplicação da Teoria dos Grafos pode ser integrada a diversas disciplinas, permitindo uma aprendizagem mais contextualizada e significativa:

\begin{itemize}
  \item \textbf{Geografia e História}: estudo de rotas comerciais, mapas, expansão territorial e o próprio contexto do Problema das Pontes de Königsberg.
  \item \textbf{Ciências e Biologia}: representação de cadeias alimentares, redes metabólicas e estudos epidemiológicos sobre a propagação de doenças.
  \item \textbf{Língua Portuguesa}: análise de redes semânticas e diagramas de organização de texto.
\end{itemize}

O cruzamento dessas áreas mostra ao estudante como a Matemática atua como uma linguagem universal para a resolução de problemas complexos que extrapolam os limites exatos da disciplina.

% ==========================================================
\chapter{PROPOSTA DIDÁTICA PARA O ENSINO DE GRAFOS}
\label{cap:proposta}
% ==========================================================

Os capítulos anteriores construíram dois pilares complementares: o
Capítulo~\ref{cap:fundamentos} estabeleceu o arcabouço matemático formal da
Teoria dos Grafos e revisou a literatura, enquanto o Capítulo~\ref{cap:bncc}
analisou as conexões com as diretrizes curriculares. O presente capítulo traduz
esses pilares em uma proposta de ensino coerente, progressiva e metodologicamente
fundamentada, apresentando atividades contextualizadas, sugestões de sequência
didática e uma discussão sobre desafios de implementação.

% ----------------------------------------------------------
\section{Metodologia de Resolução de Problemas}
% ----------------------------------------------------------

A resolução de problemas ligados ao cotidiano é uma das maneiras mais adequadas
para se abordar a Teoria dos Grafos no Ensino Médio. As Orientações Educacionais
Complementares aos Parâmetros Curriculares Nacionais (PCN+) enfatizam que a
resolução de problemas é peça central para o ensino de Matemática, pois o pensar
e o fazer se mobilizam e se desenvolvem quando o indivíduo está engajado
ativamente no enfrentamento de desafios \cite[p.~53]{brasil2002}. Essa
competência não se desenvolve quando propomos apenas exercícios de aplicação,
pois, nesse caso, o que está em ação é uma simples transposição analógica sem
garantia de transferência para situações novas \cite[p.~112]{brasil2002}.

Nessa perspectiva, cada atividade apresentada a seguir foi deliberadamente
estruturada para: (i) partir de um problema concreto e motivador; (ii) conduzir o
estudante à formalização gradual dos conceitos do
Capítulo~\ref{cap:fundamentos}; e (iii) promover a generalização e a reflexão
sobre o resultado obtido.

Adicionalmente, a LDB preconiza que as metodologias de resolução de problemas
sejam adequadas ao cotidiano discente \cite{brasil1996}, buscando fomentar o
pensamento autônomo e a criação de soluções próprias. Por isso, as atividades
utilizam contextos como redes sociais, rotas de entrega e organização de um
bairro, situações familiares aos jovens.

% ----------------------------------------------------------
\section{Atividades Sugeridas}
% ----------------------------------------------------------

\subsection{Atividade 1 — A Vila dos Grafos: construindo e analisando um grafo}
\label{ativ:vila}

\textbf{Conceitos mobilizados:}
Definição~\ref{def:grafo} (Grafo), Definição~\ref{def:grau} (Grau de um
vértice) e Lema~\ref{lema:handshaking} (Lema do Aperto de Mão).

\textbf{Objetivo:} Ao concluir esta atividade, o estudante será capaz de
representar uma situação real por meio de um grafo, identificar vértices e
arestas, calcular graus e verificar experimentalmente o Lema do Aperto de Mão.

\vspace{0.4cm}
\noindent\textbf{Situação-problema.} Você é o engenheiro de trânsito
responsável por mapear o centro de um pequeno bairro chamado ``Vila dos
Grafos''. Neste bairro, existem \textbf{5} pontos de grande importância:

\begin{itemize}[noitemsep]
    \item \textbf{A:} Prefeitura
    \item \textbf{B:} Escola
    \item \textbf{C:} Padaria
    \item \textbf{D:} Praça Central
    \item \textbf{E:} Hospital
\end{itemize}

Para conectar esses locais, existem \textbf{7 ruas de mão dupla}:

\begin{enumerate}[noitemsep]
    \item Prefeitura (A) $\leftrightarrow$ Escola (B)
    \item Prefeitura (A) $\leftrightarrow$ Praça Central (D)
    \item Escola (B) $\leftrightarrow$ Padaria (C)
    \item Escola (B) $\leftrightarrow$ Praça Central (D)
    \item Padaria (C) $\leftrightarrow$ Praça Central (D)
    \item Padaria (C) $\leftrightarrow$ Hospital (E)
    \item Praça Central (D) $\leftrightarrow$ Hospital (E)
\end{enumerate}

\vspace{0.3cm}
\noindent\textbf{Tarefas:}

\begin{enumerate}[label=\textbf{Parte \arabic*:}, leftmargin=*]
    \item Desenhe o mapa do bairro representando os locais como pontos
          (vértices) e as ruas como linhas (arestas).
    \item Liste o conjunto de vértices $\mathcal{V}$ e o conjunto de arestas
          $\mathcal{A}$. Quantos vértices e quantas arestas o grafo possui?
    \item Calcule o grau de cada vértice (ou seja, quantas ruas se conectam a
          cada local).
    \item Some todos os graus obtidos na Parte~3 e compare com o número total
          de arestas. O que você conclui? Explique a relação matemática com
          base no Lema do Aperto de Mão (Lema~\ref{lema:handshaking}).
\end{enumerate}

\vspace{0.4cm}
\noindent\textbf{Gabarito e resolução esperada.}

\vspace{0.2cm}
\noindent\textit{Parte 1 --- O grafo.}
O estudante deve desenhar um grafo equivalente ao diagrama abaixo.

\begin{center}
\begin{tikzpicture}[node distance=3cm, thick]
    \node[circle, draw, fill=gray!10, minimum size=0.9cm] (D) {D};
    \node[circle, draw, fill=gray!10, minimum size=0.9cm] (A) [above left of=D] {A};
    \node[circle, draw, fill=gray!10, minimum size=0.9cm] (B) [above right of=D] {B};
    \node[circle, draw, fill=gray!10, minimum size=0.9cm] (C) [right of=D] {C};
    \node[circle, draw, fill=gray!10, minimum size=0.9cm] (E) [below right of=D] {E};

    \draw (A) -- (B); \draw (A) -- (D); \draw (B) -- (C);
    \draw (B) -- (D); \draw (C) -- (D);
    \draw (C) -- (E);
    \draw (D) -- (E);
\end{tikzpicture}
\end{center}

\noindent\textit{Parte 2 --- Vértices e arestas.}
\begin{itemize}[noitemsep]
    \item $\mathcal{V} = \{A, B, C, D, E\}$ $\longrightarrow$ \textbf{5 vértices}
    \item $\mathcal{A} = \{AB, AD, BC, BD, CD, CE, DE\}$ $\longrightarrow$ \textbf{7 arestas}
\end{itemize}

\noindent\textit{Parte 3 --- Graus.}
\begin{itemize}[noitemsep]
    \item $\delta(A) = 2$ \quad $\delta(B) = 3$ \quad $\delta(C) = 3$
          \quad $\delta(D) = 4$ \quad $\delta(E) = 2$
\end{itemize}

\noindent\textit{Parte 4 --- Verificação do Lema do Aperto de Mão.}
\[
  \sum_{v \in \mathcal{V}} \delta(v) = 2 + 3 + 3 + 4 + 2 = 14 = 2 \times 7 = 2\,|\mathcal{A}|. \quad \checkmark
\]
Cada rua (aresta) incide em dois locais (vértices); ao somarmos os graus,
contamos cada aresta duas vezes, o que confirma o
Lema~\ref{lema:handshaking}.

% ----------------------------------------------------------
\subsection{Atividade 2 — Redes Sociais e Conectividade}
\label{ativ:redes}

\textbf{Conceitos mobilizados:}
Definição~\ref{def:grafo}, Definição~\ref{def:dirigido} (Grafo não dirigido e
dígrafo) e Definição~\ref{def:grau}.

\textbf{Objetivo:} Explorar a distinção entre grafos dirigidos e não dirigidos
a partir de um contexto familiar ao estudante: as redes sociais.

\vspace{0.4cm}
\noindent\textbf{Situação-problema.} Considere um grupo de cinco pessoas:
Ana, Bruno, Carla, Diego e Eduarda. Em uma rede de \emph{amizade mútua}
(como o Facebook), a relação é simétrica: se Ana é amiga de Bruno, então
Bruno é amigo de Ana. Em uma rede de \emph{seguidores} (como o Twitter/X),
a relação é assimétrica: Ana pode seguir Bruno sem que Bruno a siga.

As relações de \emph{amizade} existentes são:\\[4pt]
Ana--Bruno, Ana--Carla, Bruno--Diego, Carla--Diego, Diego--Eduarda, Bruno--Eduarda.

As relações de \emph{seguimento} (de $\to$ para) são:\\[4pt]
Ana $\to$ Bruno, Bruno $\to$ Diego, Carla $\to$ Ana, Carla $\to$ Diego,
Diego $\to$ Eduarda, Eduarda $\to$ Bruno.

\noindent\textbf{Tarefas:}
\begin{enumerate}[label=\textbf{Parte \arabic*:}, leftmargin=*]
    \item Represente a rede de amizades como um grafo \emph{não dirigido} e
          calcule o grau de cada vértice (pessoa).
    \item Represente a rede de seguidores como um \emph{dígrafo}. Para cada
          vértice, calcule o \emph{grau de entrada} (quantos o seguem) e o
          \emph{grau de saída} (quantos ele segue).
    \item Qual das duas representações é mais adequada para modelar uma plataforma
          de mensagens instantâneas em que dois usuários precisam se adicionar
          mutuamente para conversar? Justifique com base nas definições do
          Capítulo~\ref{cap:fundamentos}.
\end{enumerate}

% ----------------------------------------------------------
\subsection{Atividade 3 — O Problema do Caixeiro Viajante Simplificado}
\label{ativ:caixeiro}

\textbf{Conceitos mobilizados:}
Definição~\ref{def:completo} (Grafo completo $K_n$), definição de caminho
e Exemplo~\ref{ex:k4}.

\textbf{Objetivo:} Introduzir intuitivamente o conceito de otimização em
grafos, mostrando a explosão combinatória do número de rotas possíveis.

\vspace{0.4cm}
\noindent\textbf{Situação-problema.} Um entregador precisa visitar quatro
cidades --- Alfa, Beta, Gama e Delta --- partindo e retornando à cidade Alfa.
As distâncias (em km) entre cada par de cidades são dadas pela tabela:

\begin{center}
\begin{tabular}{c|cccc}
           & Alfa & Beta & Gama & Delta \\ \hline
  Alfa     &  ---  & 10   & 15   & 20    \\
  Beta     &  10   & ---  & 35   & 25    \\
  Gama     &  15   & 35   & ---  & 30    \\
  Delta    &  20   & 25   & 30   & ---   \\
\end{tabular}
\end{center}

\noindent\textbf{Tarefas:}
\begin{enumerate}[label=\textbf{Parte \arabic*:}, leftmargin=*]
    \item Represente as cidades e as distâncias como um grafo completo $K_4$
          com arestas ponderadas.
    \item Liste todas as rotas possíveis que partem de Alfa, visitam as demais
          cidades exatamente uma vez e retornam a Alfa. Quantas são?
    \item Calcule o custo (distância total) de cada rota e identifique a rota
          de menor custo.
    \item Se o número de cidades fosse 10, quantas rotas distintas existiriam?
          O que isso nos diz sobre a complexidade do problema?
\end{enumerate}

\vspace{0.3cm}
\noindent\textbf{Gabarito resumido.}

\noindent\textit{Parte 2.} As rotas são permutações de $\{B, G, D\}$:
3! = 6 rotas possíveis.

\noindent\textit{Parte 3.}
\begin{itemize}[noitemsep]
  \item A--B--G--D--A: $10+35+30+20 = 95$ km
  \item A--B--D--G--A: $10+25+30+15 = 80$ km \quad (\textbf{menor})
  \item A--G--B--D--A: $15+35+25+20 = 95$ km
  \item A--G--D--B--A: $15+30+25+10 = 80$ km \quad (igual à segunda)
  \item A--D--B--G--A: $20+25+35+15 = 95$ km
  \item A--D--G--B--A: $20+30+35+10 = 95$ km
\end{itemize}

\noindent\textit{Parte 4.} Com 10 cidades, haveria $(10-1)! = 362.880$ rotas.
Isso ilustra por que o Problema do Caixeiro Viajante é considerado um dos
grandes desafios da Computação e da Matemática Discreta.

% ----------------------------------------------------------
\section{Sequência Didática Sugerida}
% ----------------------------------------------------------

Com base na articulação entre os fundamentos matemáticos
(Capítulo~\ref{cap:fundamentos}), a análise documental (Capítulo~\ref{cap:bncc})
e as atividades propostas, sugere-se a seguinte sequência didática de
quatro encontros:

\begin{enumerate}
  \item \textbf{1\textordmasculine{} encontro --- Introdução e motivação histórica:}
        apresentação do problema das Pontes de Königsberg como motivação
        histórica e como situação-problema aberta. O professor conduz uma
        discussão coletiva antes de formalizar a Definição~\ref{def:grafo}.
  \item \textbf{2\textordmasculine{} encontro --- Construindo grafos:}
        realização da Atividade~\ref{ativ:vila} (``A Vila dos Grafos''), que
        mobiliza as Definições~\ref{def:grafo} e \ref{def:grau} e o
        Lema~\ref{lema:handshaking}. O professor sistematiza o Lema do Aperto
        de Mão ao final.
  \item \textbf{3\textordmasculine{} encontro --- Redes e sentidos:}
        discussão sobre grafos dirigidos e não dirigidos por meio da
        Atividade~\ref{ativ:redes} (``Redes Sociais e Conectividade'').
        Retoma-se o Teorema~\ref{teo:euler} como fecho da discussão sobre
        caminhos e circuitos.
  \item \textbf{4\textordmasculine{} encontro --- Otimização e síntese:}
        realização da Atividade~\ref{ativ:caixeiro} (``O Problema do Caixeiro
        Viajante Simplificado''), destacando a conexão com o pensamento
        computacional e a explosão combinatória. Encerramento com discussão
        sobre a presença dos grafos em aplicativos de navegação, redes sociais
        e outros contextos digitais.
\end{enumerate}

Recomenda-se, ao longo dos encontros, o uso de recursos digitais como o
GeoGebra \cite{silva2020}, que permite a construção interativa de grafos e
favorece a aprendizagem visual e exploratória.

% ----------------------------------------------------------
\section{Desafios e Limitações na Implementação na Rede Pública}
% ----------------------------------------------------------

A implementação de qualquer conteúdo novo na rede pública enfrenta desafios
concretos: a necessidade de formação docente específica, a sobrecarga curricular
e a escassez de recursos tecnológicos. No caso da Teoria dos Grafos, há ainda a
ausência de materiais didáticos consolidados para o Ensino Médio em língua
portuguesa.

Esta proposta, ao organizar claramente as definições
(Capítulo~\ref{cap:fundamentos}), as conexões curriculares (Capítulo~\ref{cap:bncc})
e as atividades (presentes neste capítulo), pretende minimizar esse último
obstáculo, oferecendo um ponto de partida sistematizado para professores
interessados em inovar sua prática pedagógica.

Outro desafio relevante é a carga horária limitada da Matemática no Ensino Médio.
Recomenda-se que a sequência didática proposta seja desenvolvida em projetos
interdisciplinares ou em momentos de aprofundamento, como eletivas ou clubes de
matemática, especialmente nas escolas de tempo integral. A articulação com o
pensamento computacional e a interdisciplinaridade (Seção~\ref{cap:bncc}) pode
justificar a alocação de tempo em parceria com outras disciplinas.

% ==========================================================
\chapter{CONCLUSÃO}
% ==========================================================

\textit{[O autor deve desenvolver aqui o texto conclusivo definitivo, retomando
o problema de pesquisa, sintetizando os resultados dos três capítulos centrais
e propondo encaminhamentos para pesquisas futuras.]}

% ----------------------------------------------------------
% ELEMENTOS PÓS-TEXTUAIS
% ----------------------------------------------------------
\postextual

\bibliography{referencias}

% ---
% Apêndices
% ---
\begin{apendicesenv}
\partapendices

\chapter{Subtítulo do apêndice}
Apêndices têm o objetivo de melhorar a compreensão textual, ou seja, completar
ideias desenvolvidas no decorrer do trabalho.
\end{apendicesenv}

% ---
% Anexos
% ---
\begin{anexosenv}
\partanexos

\chapter{Subtítulo do anexo}
Anexos são elementos que dão suporte ao texto, mas que não foram elaborados pelo
autor.

\end{anexosenv}
\listoftodos

\end{document}